\section{Andean (Inca)}

The \textbf{Andean} tradition in its Inca-era form survives mainly through Spanish chroniclers of the sixteenth and seventeenth centuries, whose accounts differ and were filtered through Christian categories. It belongs to the highland peoples of Peru, Bolivia, and the surrounding Andes. The central maker is \textbf{Viracocha}, also titled \textbf{Pachayachachic}, glossed as the teacher and maker of the world.

\subsection{The Creation Model}

In the beginning the world lies in darkness, with no sun and no day. The maker rises out of the great primordial lake, the wellspring of creation. This is not the first attempt. An earlier world and a first race of beings had already been shaped, but they failed and were destroyed, in many tellings by a flood that overturned that age. So the cosmos arrives through trial and erasure, the maker beginning again on the ruin of a prior world.

Beginning anew, the maker brings forth the sun, the moon, and the stars and sets their courses, driving back the dark. He fashions a new humanity from stone and clay, shaping each people in its own form and sending them underground to emerge at appointed places across the land. Creation here is a distribution of life across the ground. Having ordered the world, the maker walks among the peoples as a teacher, then departs westward across the sea toward the horizon, leaving the world running and promising return.

\subsection{Key Motifs}

The world moves through ages, each ended by a turning called \textbf{pachakuti}. The flood that ends the first creation is the water-turning, \textbf{unu pachakuti}. Each people emerges at its own origin place, a \textbf{pacarina}, so local origins coexist under the single imperial maker. The ruined city of \textbf{Tiwanaku} near \textbf{Lake Titicaca} is held as the site where the sun was set in the sky.

\begin{longtable}{@{}
ll@{}}
\toprule
Term & Sense \\
\midrule
\textbf{pacha} & world, time, space, an epoch \\
\textbf{pachakuti} & a cataclysmic turning that ends one age and starts another \\
\textbf{pacarina} & the origin place from which a people first emerged \\
\bottomrule
\end{longtable}

The Andean account fits the shared pattern of the one going out into the many and returning, since a single maker brings forth the lights and the distributed peoples and then withdraws across the sea with a promise to come back.
